\chapter{Automatic curvature control and bounded-slope uniqueness}\label{sec:graphical-global-refinements}
The local cutoff estimates of \cite{EH91}, Section 3, remove some hypotheses previously imposed to justify a maximum principle at infinity. We record exactly which bounds become conclusions. Throughout, hypersurfaces are smooth Euclidean graphs, $n\geq1$, with the conventions of Section \ref{sec:graphical-flow-angle}. Statements about arbitrary locally Lipschitz graphs with unbounded slope belong to a different existence theorem.

\begin{proposition}[Bounded slope gives automatic positive-time curvature control]\label{prop:bounded-slope-automatic-curvature}
Let $u:\mathbb R^n\times(0,T]\to\mathbb R$ be a smooth graphical flow with $v\leq V<\infty$ on the whole slab. Then
\[
 |A(x,t)|^2\leq C(n)V^4/t.
\]
In fact its sharp slope-based estimate is
\[
 (1+2t|A(x,t)|^2)v(x,t)^2\leq V^2.
\]
No curvature or spatial growth bound is assumed in either assertion. In particular they apply with $V=V_0$ to a smooth flow up to time zero with bounded initial slope, by Corollary \ref{cor:entire-slope-without-growth}.
\end{proposition}
\begin{proof}
Fix $0<\epsilon<t$. Apply Theorem \ref{thm:graphical-base-interior-curvature} on a base ball centered at the chosen point and on the time interval $[\epsilon,t]$, using time elapsed from $\epsilon$. At its center the estimate is $a(x,t)\leq C(n)V^4((t-\epsilon)^{-1}+R^{-2})$. The graph is smooth on every such compact ball. Let $R\to\infty$ and then $\epsilon\downarrow0$. This proves the first estimate, and hence uniform curvature control on every closed positive-time slab.

Theorem \ref{thm:entire-graph-curvature-smoothing} now applies on $[\epsilon,t]$: its formerly assumed slab curvature bound has just been proved. Its starting slope supremum is at most $V$. Therefore $(1+2(t-\epsilon)a(x,t))v(x,t)^2\leq V^2$. Let $\epsilon\downarrow0$ for the sharp form. No convergence of the curvature at time zero is needed. This is the local-to-global procedure discussed in \cite{EH91}, Remark 3.3(i), followed by the exact time-weighted identity already in the project.
\end{proof}

\begin{corollary}[Ancient bounded-slope graphs are affine without a curvature hypothesis]\label{cor:ancient-bounded-slope-plane}
Let $u:\mathbb R^n\times(-\infty,T)\to\mathbb R$ be a smooth entire graphical flow with one uniform bound for $v$ over all space and past times. Then $u(x,t)=p\cdot x+c$ for fixed $p,c$. The slab curvature assumption in Corollary \ref{cor:ancient-entire-graph-rigidity} is unnecessary once the local curvature estimate is available.
\end{corollary}
\begin{proof}
Fix $t<T$ and a starting time $s<t$. Translate time in Proposition \ref{prop:bounded-slope-automatic-curvature}. It gives $|A(x,t)|^2\leq C(n)V^4/(t-s)$, uniformly in $s$. Let $s\to-\infty$. Thus $A=0$ at every point and time. Lemma \ref{lem:graph-curvature-hessian-comparison} gives $D^2u=0$, so $u=p(t)\cdot x+c(t)$. The graph equation has zero right-hand side, hence $p'=0$ and $c'=0$. An entire translator with a bounded slope is consequently a stationary affine graph; translators on strips or with unbounded slope are outside these assumptions.
\end{proof}

\begin{corollary}[A bounded-gradient entire minimal graph is affine in every dimension]\label{cor:bounded-gradient-minimal-graph-affine}
If a smooth solution $u:\mathbb R^n\to\mathbb R$ of the minimal graph equation has $\sup|Du|<\infty$, then $u$ is affine, for every $n\geq1$. This is a bounded-gradient conclusion, not the unrestricted Bernstein theorem. The stationary use of graphical-flow estimates is explained in \cite{EH91}, Remark 3.6(iii).
\end{corollary}
\begin{proof}
Regard the graph as the time-independent solution $u(x,t)=u(x)$ for all $t\in\mathbb R$. The graphical flow equation is satisfied because its mean curvature is zero. Its slope bound is uniform for every past time. Corollary \ref{cor:ancient-bounded-slope-plane} therefore gives the affine conclusion. The argument leaves open the unrestricted entire minimal-graph problem when the slope is unbounded; no dimension threshold is being removed from that different statement.
\end{proof}

\begin{lemma}[A maximum principle with a time-integrable drift bound]\label{lem:maximum-integrable-time-drift}
Suppose $0\leq a(x,t)\leq A_0I$ and $|b(x,t)|\leq B(t)$ on $\mathbb R^n\times(0,T]$, where $B$ is continuous, nonnegative and integrable on $(0,T]$. Let $f$ be continuous on $\mathbb R^n\times[0,T]$, smooth enough for the following inequality at positive time, and satisfy
\[
 f_t-a^{ij}f_{ij}-b^if_i\leq0,\quad f(x,0)\leq M,
 \quad f(x,t)\leq C(1+|x|^2)^m
\]
uniformly on the slab for some $C,m\geq0$. Then $f\leq M$. Continuity at time zero is local; a uniform spatial rate of initial convergence is not assumed.
\end{lemma}
\begin{proof}
Repeat the barrier proof of Lemma \ref{lem:entire-polynomial-maximum-principle} with $q>\max\{m,1\}$ and
\[
 E(x,t)=(1+|x|^2)^q
 \exp\left((1+4nA_0q^2)t+q\int_0^t B(s)\,ds\right).
\]
The spatial derivative bounds in that lemma give $E_t-a^{ij}E_{ij}-b^iE_i\geq E>0$ for $t>0$. Integrability of $B$ makes $E$ finite and continuous at time zero. For fixed $\eta>0$, the uniform polynomial upper bound confines every possible positive maximum of $f-M-\eta E$ to a fixed compact base ball, uniformly in time. Initial values there are strictly negative. Continuity and the strict parabolic inequality exclude a positive maximum at later time exactly as before. Let $\eta\downarrow0$ at each point. The argument allows $B(t)=Ct^{-1/2}$, but not a nonintegrable coefficient merely by ignoring its singularity.
\end{proof}

\begin{theorem}[Comparison and uniqueness in the bounded-slope entire class]\label{thm:entire-bounded-slope-comparison}
Let $u_1,u_2$ be smooth entire graphical flows for $0<t\leq T$, continuous locally up to $t=0$, with one finite bound for both slopes on the whole positive-time slab. If $u_1(x,0)\leq u_2(x,0)$ for every $x$, then $u_1\leq u_2$ for every time. Thus equal initial functions give equal flows in this class, without an assumed initial or slab curvature bound. The curvature input is the locally proved estimate based on \cite{EH91}, Theorem 3.1.
\end{theorem}
\begin{proof}
Put $w=u_1-u_2$. Subtract the equations in the same way as Proposition \ref{prop:entire-graph-bounded-geometry-uniqueness}. Then
\[
 w_t=a^{ij}(Du_1)w_{ij}+b^kw_k,
 \qquad
 b^k=\int_0^1\partial_{p_k}a^{ij}(Du_2+\theta Dw)\,d\theta\,(u_2)_{ij}.
\]
The slope bounds control derivatives of the coefficient matrix on every interpolated gradient. Proposition \ref{prop:bounded-slope-automatic-curvature} and Lemma \ref{lem:graph-curvature-hessian-comparison} give $|D^2u_2|\leq C t^{-1/2}$, so $|b|\leq C t^{-1/2}$. This is integrable in time. The slope bounds and continuity of the two heights at the base origin give $|w(x,t)|\leq C_T(1+|x|)$ uniformly for $0\leq t\leq T$. Apply Lemma \ref{lem:maximum-integrable-time-drift} to $w$, whose initial value is nonpositive. This proves comparison. Applying it in both orders proves uniqueness. No differentiability of the initial function enters this argument.
\end{proof}

\begin{corollary}[Uniqueness of the globally Lipschitz-data construction]\label{cor:globally-lipschitz-graph-wellposedness}
For every globally Lipschitz entire initial function, the Ecker--Huisken solution from Theorem \ref{thm:ecker-huisken-global-graph-existence} is unique among smooth positive-time entire graphical flows converging locally continuously to that function and having bounded slope on each finite positive-time slab, uniformly down to time zero. Ordered initial functions produce ordered solutions.
\end{corollary}
\begin{proof}
The cited existence theorem supplies such a solution and its slope bound. On any finite time interval, another flow in the stated class satisfies the hypotheses of Theorem \ref{thm:entire-bounded-slope-comparison}. Equal initial values give equality on that interval. Since the interval is arbitrary, uniqueness holds for all time. The same theorem applied to two constructed solutions proves order preservation. Existence remains the external parabolic theorem; the comparison and uniqueness argument is local to this section. This does not assert uniqueness for unrestricted entire graphs with unbounded slopes.
\end{proof}

\begin{proposition}[Convex initial data give upward motion]\label{prop:convex-graph-data-upward-motion}
Let $f:\mathbb R^n\to\mathbb R$ be convex and globally Lipschitz, and let $u$ be its solution from Corollary \ref{cor:globally-lipschitz-graph-wellposedness}. Then
\[
 u(x,t)\geq f(x),\qquad u(x,t+s)\geq u(x,t)\quad(s,t\geq0),
 \qquad u_t\geq0\quad(t>0).
\]
If $f$ is not affine, then $u_t>0$ everywhere at positive time, and $u(x,t)>f(x)$. Thus the scalar mean curvature with the upward orientation satisfies $H<0$. This assertion concerns the sign of motion and mean curvature; it does not establish preservation of convexity of the whole second fundamental form.
\end{proposition}
\begin{proof}
A finite convex function on $\mathbb R^n$ is the supremum of its supporting affine functions, with a supporting function attaining the value at every point. We use this finite-dimensional supporting-hyperplane fact as convex-analysis background. If $f$ is $L$-Lipschitz, the slopes of these supporting functions have norm at most $L$: testing their defining inequality along each unit direction gives that bound. Each affine function defines a stationary bounded-slope graphical flow. Theorem \ref{thm:entire-bounded-slope-comparison} compares it with $u$, proving $u(\cdot,t)\geq f$ after taking the supremum over supports.

Fix $s>0$. The flow $u(x,t+s)$ has initial value $u(x,s)\geq f(x)$, remains bounded in slope on finite slabs and is smooth at $t=0$. Comparison with $u(x,t)$ proves time monotonicity. Smoothness then gives $q=u_t\geq0$ for positive time. Differentiating the graphical equation gives
\[
 q_t=a^{ij}(Du)q_{ij}
       +\partial_{p_k}a^{ij}(Du)u_{ij}q_k.
\]
On every positive-time compact cylinder this is a smooth uniformly parabolic linear equation. If $q$ vanished at an interior point at time $t_0>0$, the classical local strong maximum principle, propagated through overlapping bounded cylinders, would give $q=0$ throughout $\mathbb R^n\times(0,t_0]$. Thus $u$ would be stationary on that interval, and its locally uniform initial trace $f$ would equal this smooth stationary graph. Since $f$ is convex, $D^2f$ is positive semidefinite; the minimal graph equation $a^{ij}(Df)f_{ij}=0$ with positive definite $a$ forces $D^2f=0$. Hence $f$ would be affine. For nonaffine $f$ this contradiction proves $u_t>0$. Integrating on any closed positive-time subinterval and using $u\geq f$ gives $u(x,t)>f(x)$. Finally $u_t=-Hv$ with $v>0$ gives $H<0$. The external input for existence remains \cite{EH89}, Theorem 4.6; the local strong maximum principle and convex supporting-hyperplane fact are the background used here.
\end{proof}
